\documentclass[11pt]{article}
\usepackage[a4paper,margin=27mm]{geometry}
\usepackage{amsmath,amssymb,mathtools}
\usepackage{fontspec,unicode-math}
\setmainfont{Latin Modern Roman}
\setmathfont{Latin Modern Math}
\usepackage{xurl}
\usepackage[hidelinks]{hyperref}
\newcommand{\sref}[2]{\hyperlink{source:#1:#2}{[S#2]}}
\newcommand{\cataloguescope}[1]{\paragraph{Mathematical question.}#1}
\hypersetup{pdftitle={Problem statement}}
\begin{document}
\section{Bounded prime-sum criterion for rational points}\label{statement}

\subsection{Definitions and mathematical statement}

\paragraph{Shared notation.}$\mathbb N=\{1,2,\ldots\}$, and $\mathbb Z,\mathbb Q,\mathbb R,\mathbb C$ denote the integers, rationals, reals and complex numbers. Any other conventions are specified locally.


Let $E/\mathbb Q$ be an elliptic curve, represented by an integral equation $y^2=x^3+Ax+B$ with $A,B\in\mathbb Z$ and $4A^3+27B^2\ne0$. For every prime $p$, define
\[
N_p=1+\#\{(x,y)\in\mathbb F_p^2:y^2=x^3+Ax+B\},\qquad
f_E(X)=\sum_{\substack{p\text{ prime}\\p\le X}}\log\!\left(\frac{N_p}{p}\right)\quad(X\ge2).
\]
The extra one accounts for the point at infinity: at good primes $N_p=\#E(\mathbb F_p)$, rather than this projective count plus one. At the finitely many bad primes the same displayed affine-count formula specifies $N_p$. A change to another integral model affects at most finitely many local counts and therefore changes $f_E(X)$ by a constant for all sufficiently large $X$. Bounded means that $|f_E(X)|\le C_E$ for all real $X\ge2$, for some finite constant $C_E$.
\paragraph{Statement recorded in the supplied source.}
For every elliptic curve $E/\mathbb Q$, is $f_E(X)$ bounded as $X\to\infty$ if and only if the group $E(\mathbb Q)$ is finite? \sref{statement}{2}
\subsection{Short English statement}
Does boundedness of the logarithmic prime sum of local point counts exactly characterize elliptic curves with only finitely many rational points?
\subsection{Sources}
\begin{description}
\item[\hypertarget{source:statement:1}{S1}] ULAM.AI and the UnsolvedMath Contributors, UnsolvedMath: A Curated Collection of Open Mathematics Problems, catalogue entry. CC BY 4.0. \url{https://huggingface.co/datasets/ulamai/UnsolvedMath}.
\item[\hypertarget{source:statement:2}{S2}] Christian Wuthrich, \emph{Open problems in number theory}, lecture slides (December 2011) slides beginning Let $N_p$ be the number of solutions of $E$ modulo $p$ plus $1$, followed by Birch and Swinnerton-Dyer conjecture. \url{https://www.maths.nottingham.ac.uk/plp/pmzcw/download/open.pdf}.
\end{description}
\label{end:statement}
\end{document}
